% ==================== CHƯƠNG 1: TASK 1 (HOÀN CHỈNH - TUYẾT TRÂM NÂNG CẤP) ====================
\section{Khai thác và Đặc tả Yêu cầu (Task 1)}

\subsection{Bối cảnh nghiệp vụ, Stakeholder và Phạm vi dự án}
\textit{(Phụ trách thực hiện: Phan Lê Ngọc Như Ý)}

\subsubsection{Bối cảnh nghiệp vụ}
BrewLite là hệ thống đặt món và quản lý hoạt động của chuỗi cửa hàng đồ uống. Hệ thống hướng đến việc chuẩn hóa quy trình từ xem thực đơn, lựa chọn sản phẩm, đặt món, thanh toán không tiền mặt đến theo dõi trạng thái đơn hàng trên một nền tảng web tối ưu cho thiết bị di động (Mobile-first/PWA). Đồng thời, hệ thống hỗ trợ cửa hàng quản lý tập trung thông tin sản phẩm, tồn kho, đơn hàng, khuyến mãi và giao dịch điểm thưởng.

Trong quy trình rời rạc, khách hàng có thể phải tiếp cận thực đơn qua nhiều kênh khác nhau, khó nắm bắt tình trạng còn hàng và trạng thái xử lý đơn. Nhân viên tiếp nhận yêu cầu thủ công có thể nhầm lẫn về sản phẩm, số lượng, Size, Topping hoặc ghi chú. Việc quản lý giá bán, tồn kho, thanh toán và điểm thưởng ở nhiều nơi cũng làm tăng nguy cơ sai lệch dữ liệu và khó đối soát.

BrewLite giải quyết các vấn đề trên bằng một luồng nghiệp vụ thống nhất. Khách hàng có thể duyệt thực đơn, tìm kiếm và lọc sản phẩm, xem thông tin chi tiết, lựa chọn Size/Topping, quản lý giỏ hàng, xác nhận đơn, thanh toán và theo dõi trạng thái. Các nghiệp vụ quan trọng như kiểm tra tồn kho, tính tổng tiền, áp dụng khuyến mãi và ghi nhận điểm thưởng được xử lý tại backend nhằm hạn chế sai lệch giữa giao diện và dữ liệu hệ thống.

Các mục tiêu nghiệp vụ chính gồm:
\begin{enumerate}
    \item Cung cấp thực đơn tập trung, hỗ trợ tìm kiếm, lọc danh mục và xem chi tiết sản phẩm.
    \item Hạn chế sai sót khi khách hàng lựa chọn sản phẩm, Size, Topping và số lượng.
    \item Chuẩn hóa quy trình tạo đơn, kiểm tra tồn kho và tính tiền tại máy chủ.
    \item Hỗ trợ thanh toán không tiền mặt, khuyến mãi và tích điểm có kiểm soát.
    \item Minh bạch trạng thái đơn hàng để khách hàng và nhân viên cùng theo dõi.
    \item Xây dựng nền tảng có thể kiểm thử, triển khai và mở rộng theo kiến trúc của dự án.
\end{enumerate}

\subsubsection{Các bên liên quan (Stakeholders)}
\begin{table}[H]
\centering
\small
\renewcommand{\arraystretch}{1.25}
\caption{Các bên liên quan của hệ thống BrewLite}
\begin{tabularx}{\textwidth}{|p{0.23\textwidth}|p{0.25\textwidth}|X|}
\hline
\textbf{Stakeholder} & \textbf{Vai trò} & \textbf{Nhu cầu và Kỳ vọng chính} \\
\hline
Khách hàng (Customer) & Người sử dụng trực tiếp & Xem thực đơn, tìm kiếm/lọc món, chọn Size/Topping, đặt và thanh toán không tiền mặt, theo dõi đơn và điểm thưởng. \\
\hline
Nhân viên quầy & Tiếp nhận và hỗ trợ đơn & Tra cứu sản phẩm, xem thông tin đơn đã thanh toán (\texttt{PAID}), áp dụng ưu đãi nội bộ (\texttt{STAFF20}) và cập nhật trạng thái theo quyền được cấp. \\
\hline
Nhân viên pha chế & Thực hiện món & Nhận danh sách món, số lượng, Size/Topping và ghi chú chính xác. \\
\hline
Quản lý cửa hàng & Theo dõi vận hành & Quản lý sản phẩm, giá, tồn kho (\texttt{stockQuantity}), đơn hàng và chương trình khuyến mãi. \\
\hline
Quản trị hệ thống & Quản trị tài khoản và cấu hình & Quản lý người dùng/vai trò, kiểm soát truy cập và theo dõi tình trạng hệ thống. \\
\hline
Đơn vị cung cấp thanh toán & Hệ thống bên ngoài & Tiếp nhận yêu cầu thanh toán hợp lệ (Ví MoMo QR / Thẻ) và trả kết quả giao dịch để đối soát, đảm bảo tính bất biến (Idempotency). \\
\hline
Nhóm phát triển (6 SV) \& Giảng viên & Xây dựng và nghiệm thu sản phẩm & Yêu cầu, sơ đồ UML, kiến trúc, mã nguồn, kiểm thử tự động, CI/CD và tài liệu phải thống nhất, truy vết được. \\
\hline
\end{tabularx}
\end{table}

\subsubsection{Phạm vi dự án (In-scope \& Out-of-scope)}
\begin{itemize}
    \item \textbf{Trong phạm vi (In-scope):}
    \begin{itemize}
        \item \textit{Thực đơn và sản phẩm:} danh sách sản phẩm, tìm kiếm, lọc theo danh mục, phân trang, xem chi tiết, giá, mô tả, hình ảnh và tình trạng còn hàng; lựa chọn Size/Topping trên giao diện.
        \item \textit{Tài khoản:} đăng ký, đăng nhập, xác thực JWT và xem hồ sơ người dùng.
        \item \textit{Giỏ hàng và đặt món:} thêm, chỉnh sửa, xóa món; thay đổi số lượng và tùy chọn; xác nhận đơn; kiểm tra tồn kho và tính tổng tiền ở backend.
        \item \textit{Thanh toán và ưu đãi:} thanh toán Ví MoMo (tạo mã QR thanh toán), Thẻ tín dụng/ghi nợ mô phỏng, khóa chống trùng lặp \texttt{Idempotency-Key}, áp dụng mã khuyến mãi (\texttt{BREWLITE10}, \texttt{STAFF20}, \texttt{SGU2026}, \texttt{FREESHIP}) và ghi nhận điểm thưởng Loyalty.
        \item \textit{Theo dõi đơn hàng:} xem xác nhận đơn vừa đặt, lịch sử đơn (\texttt{GET /orders/me}) và tiến trình xử lý đơn.
        \item \textit{Dữ liệu và vận hành:} lưu trữ 7 thực thể trong PostgreSQL qua Prisma; endpoint kiểm tra sức khỏe \texttt{GET /health}; kiểm thử tự động 3 tầng và triển khai CI/CD.
    \end{itemize}
    \item \textbf{Ngoài phạm vi (Out-of-scope):}
    \begin{itemize}
        \item Ứng dụng native riêng cho iOS/Android; MVP tập trung vào web responsive theo định hướng PWA.
        \item Tích hợp thiết bị POS, máy in hóa đơn, máy quét mã vạch hoặc phần cứng chuyên dụng tại quầy.
        \item Quản lý mua hàng, nhà cung cấp, công nợ, kế toán tổng hợp, nhân sự và tính lương.
        \item Đặt bàn, quản lý chỗ ngồi, điều phối giao hàng hoặc theo dõi lộ trình shipper.
        \item Kết nối trực tiếp cổng ngân hàng thương mại thực tế (sử dụng dịch vụ thanh toán sandbox/mô phỏng).
        \item Dự báo nhu cầu hay gợi ý món bằng AI.
    \end{itemize}
\end{itemize}

\subsubsection{Giả định, ràng buộc và rủi ro}
\begin{enumerate}
    \item Giá sản phẩm được lưu dưới dạng số nguyên theo đơn vị đồng Việt Nam (VND).
    \item Dữ liệu sản phẩm, tài khoản mẫu và mã khuyến mãi được khởi tạo thông qua cơ chế Database Seed chạy lặp an toàn.
    \item Phụ thu Size/Topping được chốt thống nhất và tính toán lại ở backend khi tạo đơn hàng (\texttt{POST /orders}) để đảm bảo tính toàn vẹn dữ liệu.
    \item MVP sử dụng số lượng tồn kho chung (\texttt{stockQuantity}) cho từng sản phẩm trên toàn hệ thống.
    \item PWA yêu cầu chạy trên HTTPS (Production) hoặc \texttt{localhost} (môi trường phát triển) để kích hoạt Service Worker.
    \item Hệ thống sử dụng Next.js, NestJS, PostgreSQL và Prisma trong cấu trúc Monorepo; các thành viên tích hợp thông qua Git branch và Pull Request.
\end{enumerate}

\begin{table}[H]
\centering
\small
\renewcommand{\arraystretch}{1.25}
\caption{Rủi ro chính của dự án và hướng giảm thiểu}
\begin{tabularx}{\textwidth}{|p{0.27\textwidth}|p{0.24\textwidth}|X|}
\hline
\textbf{Rủi ro} & \textbf{Ảnh hưởng} & \textbf{Hướng giảm thiểu} \\
\hline
Nhiều yêu cầu đặt cùng sản phẩm khi tồn kho thấp & Bán vượt tồn kho (tồn kho âm) & Kiểm tra và cập nhật tồn kho trong cùng Database Transaction tại backend; kiểm thử đồng thời. \\
\hline
Giá Size/Topping không đồng nhất giữa UI và backend & Sai tổng tiền đơn hàng & Thống nhất bảng phụ thu; backend là nguồn tính tiền cuối cùng. \\
\hline
Yêu cầu thanh toán bị gửi lặp lại do mạng lỗi & Trừ tiền hoặc cộng điểm thưởng nhiều lần & Dùng header \texttt{Idempotency-Key} duy nhất kết hợp Transaction. \\
\hline
Các nhánh thành viên thay đổi schema/API không đồng bộ & Lỗi tích hợp, xung đột merge & Thống nhất API/Prisma contract trước khi code và kiểm tra qua Pull Request. \\
\hline
Lộ secret hoặc dữ liệu mật khẩu người dùng & Mất an toàn bảo mật & Không commit file \texttt{.env}; dùng biến môi trường và băm mật khẩu bằng \texttt{bcrypt}. \\
\hline
Kết nối mạng không ổn định khi khách xem thực đơn & Gián đoạn trải nghiệm & Hiển thị rõ 3 trạng thái Loading/Empty/Error, nút tải lại và cache tĩnh PWA. \\
\hline
\end{tabularx}
\end{table}

% ==================== MỤC 1.2 ====================
\subsection{Đặc tả Yêu cầu Chức năng (FR) và Phi chức năng (NFR)}
\textit{(Phụ trách thực hiện: Nguyễn Thị Tuyết Trâm phối hợp cùng Trần Thị Trúc Đào)}

\subsubsection{Danh mục Yêu cầu chức năng (Functional Requirements)}
Hệ thống BrewLite gồm 20 yêu cầu chức năng được phân thành 7 nhóm nghiệp vụ chính:

\small
\renewcommand{\arraystretch}{1.22}
\begin{longtable}{|p{0.09\textwidth}|p{0.22\textwidth}|p{0.10\textwidth}|p{0.51\textwidth}|}
\caption{Danh mục Yêu cầu chức năng (FR-01 đến FR-20)} \\
\hline
\textbf{Mã} & \textbf{Yêu cầu} & \textbf{MoSCoW} & \textbf{Tiêu chí chấp nhận (Acceptance Criteria)} \\
\hline
\endfirsthead
\hline
\textbf{Mã} & \textbf{Yêu cầu} & \textbf{MoSCoW} & \textbf{Tiêu chí chấp nhận (tiếp theo)} \\
\hline
\endhead
\hline
\endfoot
\multicolumn{4}{|l|}{\textbf{NHÓM 1: THỰC ĐƠN (MENU)}} \\
\hline
FR-01 & Xem thực đơn & Must & Mở trang thực đơn hiển thị danh sách đồ uống kèm tên, giá, tình trạng tồn kho (\texttt{GET /products}); hỗ trợ tìm kiếm, lọc danh mục, phân trang; xử lý đủ trạng thái Loading, Empty, Error; món hết hàng hiện nhãn ``Hết hàng''. \\
\hline
FR-02 & Xem chi tiết món & Must & Chọn món hiển thị đúng mô tả, giá gốc và tùy chọn (\texttt{GET /products/:id}); truy cập ID không tồn tại hiện thông báo ``Không tìm thấy món''. \\
\hline
FR-03 & Chọn size, topping \& tính đơn giá & Must & Bắt buộc chọn 1 Size (S: +0đ, M: +5.000đ, L: +10.000đ). Topping tùy chọn (Trân châu +10.000đ, Kem +10.000đ, Thạch +8.000đ), chọn được nhiều loại. Đơn giá cập nhật ngay khi đổi lựa chọn. \\
\hline
\multicolumn{4}{|l|}{\textbf{NHÓM 2: GIỎ HÀNG (CART)}} \\
\hline
FR-04 & Thêm món vào giỏ & Must & Thêm món hợp lệ tạo 1 dòng trong giỏ; trùng món và trùng tùy chọn thì cộng dồn số lượng; khác tùy chọn tạo dòng mới. \\
\hline
FR-05 & Sửa số lượng, xóa món & Must & Tăng/giảm số lượng (trong khoảng 1--20) và xóa món; giảm về 0 tự động xóa dòng khỏi giỏ. \\
\hline
FR-06 & Tính tổng tiền giỏ hàng & Must & Thành tiền từng dòng = đơn giá $\times$ số lượng; tổng tạm tính = tổng các thành tiền, sai số 0đ và cập nhật tức thì. \\
\hline
FR-07 & Lưu trữ giỏ hàng & Should & Tải lại trang hoặc đóng mở lại trình duyệt vẫn giữ nguyên trạng thái giỏ hàng. \\
\hline
\multicolumn{4}{|l|}{\textbf{NHÓM 3: XÁC THỰC (AUTHENTICATION)}} \\
\hline
FR-08 & Đăng ký tài khoản & Must & Tạo tài khoản với email và mật khẩu hợp lệ (\texttt{POST /auth/register}); từ chối email trùng hoặc mật khẩu dưới 6 ký tự / thiếu chữ và số; băm mật khẩu trước khi lưu. \\
\hline
FR-09 & Đăng nhập & Must & Đăng nhập thành công cấp JWT token (\texttt{POST /auth/login}); nhập sai báo lỗi chung ``Email hoặc mật khẩu không đúng''. \\
\hline
FR-10 & Bảo vệ chức năng cần đăng nhập & Must & Khách chưa đăng nhập vẫn xem được menu và giỏ hàng, nhưng khi đặt đơn, thanh toán hoặc xem lịch sử đơn sẽ bị từ chối (401) và chuyển hướng đăng nhập. \\
\hline
\multicolumn{4}{|l|}{\textbf{NHÓM 4: ĐẶT ĐƠN (ORDER)}} \\
\hline
FR-11 & Tạo đơn hàng & Must & Khách đã đăng nhập với giỏ không rỗng tạo được đơn trạng thái \texttt{PENDING} (\texttt{POST /orders}); chốt giá và khoản giảm tại thời điểm đặt. \\
\hline
FR-12 & Kiểm tra và trừ tồn kho & Must & Kiểm tra \texttt{stockQuantity >= quantity} trong Database Transaction; không đủ hàng thì rollback và báo rõ món thiếu; đặt thành công trừ đúng tồn kho, không bao giờ âm. \\
\hline
FR-13 & Áp mã khuyến mãi & Should & Chấp nhận các mã hợp lệ (\texttt{BREWLITE10}, \texttt{STAFF20}, \texttt{SGU2026}, \texttt{FREESHIP}); kiểm tra quyền nhân viên (\texttt{staffOnly}) và giới hạn lượt dùng; tổng sau giảm làm tròn xuống VND nguyên và không âm. \\
\hline
\multicolumn{4}{|l|}{\textbf{NHÓM 5: THANH TOÁN (PAYMENT)}} \\
\hline
FR-14 & Thanh toán đơn qua cổng mô phỏng & Must & Thanh toán đơn \texttt{PENDING} qua Ví MoMo (tạo mã QR) / Thẻ (\texttt{POST /payments}); thành công chuyển đơn sang \texttt{PAID}; đơn đã \texttt{PAID} không cho thanh toán lại. \\
\hline
FR-15 & Chống thanh toán trùng (Idempotency) & Must & Yêu cầu thiếu \texttt{Idempotency-Key} bị từ chối (400); gửi trùng khóa trả về kết quả cũ, không tạo giao dịch mới và không cộng điểm lần 2. \\
\hline
FR-16 & Xử lý thanh toán thất bại & Must & Khi nhận kết quả \texttt{PAYMENT\_FAILED}, đơn giữ nguyên \texttt{PENDING}, giữ nguyên giỏ hàng và cho phép khách thử thanh toán lại. \\
\hline
\multicolumn{4}{|l|}{\textbf{NHÓM 6: XÁC NHẬN VÀ LỊCH SỬ ĐƠN}} \\
\hline
FR-17 & Màn hình xác nhận đơn & Must & Hiển thị mã đơn, danh sách món, tổng tiền, trạng thái sau khi thanh toán thành công và tự động làm rỗng giỏ hàng. \\
\hline
FR-18 & Lịch sử đơn của tôi & Must & Xem danh sách đơn của chính mình xếp theo mới nhất (\texttt{GET /orders/me}); ngăn chặn truy cập đơn của người dùng khác. \\
\hline
\multicolumn{4}{|l|}{\textbf{NHÓM 7: KHÁCH HÀNG THÂN THIẾT \& HỒ SƠ}} \\
\hline
FR-19 & Tích điểm thưởng (Loyalty) & Should & Đơn chuyển \texttt{PAID} cộng 1 điểm cho mỗi 10.000đ thực trả (làm tròn xuống); mỗi đơn chỉ cộng điểm 1 lần duy nhất. \\
\hline
FR-20 & Xem thông tin cá nhân & Should & Hiển thị họ tên, email, tổng điểm tích lũy (\texttt{GET /users/me}) và trạng thái các đơn hàng đã đặt. \\
\hline
\end{longtable}
\normalsize

\subsubsection{Danh mục Yêu cầu phi chức năng (Non-Functional Requirements)}
\small
\renewcommand{\arraystretch}{1.22}
\begin{longtable}{|p{0.09\textwidth}|p{0.16\textwidth}|p{0.22\textwidth}|p{0.10\textwidth}|p{0.35\textwidth}|}
\caption{Danh mục Yêu cầu phi chức năng (NFR-01 đến NFR-11)} \\
\hline
\textbf{Mã} & \textbf{Nhóm} & \textbf{Yêu cầu} & \textbf{MoSCoW} & \textbf{Tiêu chí chấp nhận đo kiểm được} \\
\hline
\endfirsthead
\hline
\textbf{Mã} & \textbf{Nhóm} & \textbf{Yêu cầu} & \textbf{MoSCoW} & \textbf{Tiêu chí chấp nhận (tiếp theo)} \\
\hline
\endhead
\hline
\endfoot
NFR-01 & Hiệu năng & Thời gian phản hồi API & Must & Các thao tác đọc menu, giỏ hàng và tạo đơn phản hồi $\le 500\text{ms}$ ở tải thông thường. \\
\hline
NFR-02 & Bảo mật & Bảo vệ mật khẩu & Must & Mật khẩu băm bằng \texttt{bcrypt}; không lưu mật khẩu dạng rõ trong CSDL hay log. \\
\hline
NFR-03 & Bảo mật & Mã hóa đường truyền & Must & Môi trường Production chạy qua HTTPS (TLS $\ge 1.2$). \\
\hline
NFR-04 & Bảo mật & Kiểm soát truy cập Server & Must & Kiểm tra quyền sở hữu phía server qua JWT; truy cập đơn người khác trả về 403/404. \\
\hline
NFR-05 & Khả dụng & Giao diện Responsive & Must & Hiển thị và thao tác chuẩn từ 360px (mobile) đến 1920px (desktop), không cuộn ngang. \\
\hline
NFR-06 & Khả dụng & Thông báo lỗi rõ ràng & Should & Hiển thị lỗi tiếng Việt dễ hiểu, không làm lộ stack trace kỹ thuật ra UI. \\
\hline
NFR-07 & Triển khai & Docker \& Health Check & Must & Chạy \texttt{docker compose up} khởi động được toàn bộ hệ thống; endpoint \texttt{GET /health} trả về 200 OK. \\
\hline
NFR-08 & Triển khai & Biến môi trường & Must & Secret, JWT key, chuỗi kết nối DB quản lý qua file \texttt{.env} và cấu hình Cloud, không hard-code. \\
\hline
NFR-09 & Triển khai & Pipeline CI tự động & Should & Mỗi Pull Request tự động chạy Lint, Test, Build; pipeline đỏ chặn merge. \\
\hline
NFR-10 & Triển khai & Tài liệu triển khai & Must & README và báo cáo hướng dẫn chi tiết chạy local, test, deploy Cloud/Vercel, rollback. \\
\hline
NFR-11 & Dữ liệu & Toàn vẹn giao dịch & Must & Dùng Database Transaction đảm bảo không có đơn \texttt{PAID} thiếu bản ghi thanh toán; tiền lưu số nguyên VND. \\
\hline
\end{longtable}
\normalsize

% ==================== MỤC 1.3 ====================
\subsection{Sơ đồ Use-case tổng quát và Đặc tả Use-case trọng tâm}
\textit{(Phụ trách thực hiện: Nguyễn Thị Tuyết Trâm \& Phan Lê Ngọc Như Ý)}

\subsubsection{Sơ đồ Use-case tổng quát}
\begin{figure}[H]
\centering
\chenHinh{width=0.95\textwidth,height=0.58\textheight,keepaspectratio}{../diagrams/usecase-overview.png}
\caption{Sơ đồ Use-case tổng quát của hệ thống đặt đồ uống BrewLite}
\end{figure}

\subsubsection{Đặc tả Use-case trọng tâm: Đặt đơn và Thanh toán}
\begin{figure}[H]
\centering
\chenHinh{width=0.95\textwidth,height=0.52\textheight,keepaspectratio}{../diagrams/usecase-checkout.png}
\caption{Sơ đồ Use-case chi tiết cho nghiệp vụ Đặt đơn và Thanh toán}
\end{figure}

\begin{itemize}
    \item \textbf{Actor chính:} Khách hàng (\textit{Customer}). \textbf{Actor phụ:} Cổng thanh toán mô phỏng (\textit{Payment Gateway System}).
    \item \textbf{Liên kết yêu cầu:} \texttt{FR-10} đến \texttt{FR-17}, \texttt{FR-19}.
    \item \textbf{Tiền điều kiện:} Giỏ hàng có ít nhất 1 sản phẩm; khách hàng đã đăng nhập với JWT token hợp lệ.
    \item \textbf{Hậu điều kiện (Thành công):} Đơn hàng được tạo và chuyển trạng thái \texttt{PAID}; tồn kho sản phẩm giảm đúng số lượng đặt; lưu giao dịch thanh toán gắn \texttt{Idempotency-Key}; cộng điểm thưởng vào tài khoản; xóa giỏ hàng và hiển thị màn hình xác nhận đơn.
    \item \textbf{Hậu điều kiện (Thất bại):} Nếu chưa tạo đơn (lỗi tồn kho/chưa đăng nhập) thì không phát sinh đơn mới và giữ nguyên tồn kho. Nếu thanh toán thất bại (\texttt{PAYMENT\_FAILED}), đơn giữ trạng thái \texttt{PENDING}, giữ nguyên giỏ hàng và không cộng điểm thưởng để khách có thể thanh toán lại.
    \item \textbf{Luồng chính (Main Flow):}
    \begin{enumerate}
        \item Khách hàng nhấn ``Thanh toán'' tại màn hình giỏ hàng; hệ thống hiển thị danh sách món và tổng tạm tính (\texttt{UC09}).
        \item Hệ thống kiểm tra trạng thái đăng nhập (\texttt{UC08}).
        \item Khách hàng có thể nhập mã khuyến mãi (\texttt{BREWLITE10}, \texttt{STAFF20}...); hệ thống kiểm tra tính hợp lệ và cập nhật tổng tiền sau giảm giá (\texttt{UC10}, \texttt{UC10a}).
        \item Khách hàng chọn phương thức thanh toán không tiền mặt (Ví MoMo QR / Thẻ) và xác nhận (\texttt{UC11}).
        \item Hệ thống gọi \texttt{POST /orders} tạo đơn hàng trạng thái \texttt{PENDING}, kiểm tra và trừ tồn kho trong Database Transaction.
        \item Hệ thống sinh mã \texttt{Idempotency-Key} và gọi \texttt{POST /payments} gửi yêu cầu sang Cổng thanh toán.
        \item Cổng thanh toán trả về thành công; hệ thống cập nhật trạng thái đơn từ \texttt{PENDING} sang \texttt{PAID}.
        \item Hệ thống cộng điểm thưởng vào tài khoản khách hàng (\texttt{UC12}).
        \item Hệ thống xóa giỏ hàng và hiển thị màn hình xác nhận đơn hàng (\texttt{UC13}).
        \item Khách hàng theo dõi trạng thái xử lý đơn hàng (\texttt{UC17}).
    \end{enumerate}
    \item \textbf{Luồng thay thế (Alternative Flows):}
    \begin{itemize}
        \item \textbf{A1 -- Không dùng mã khuyến mãi:} Bỏ qua bước 3, giữ nguyên tổng tạm tính.
        \item \textbf{A2 -- Đổi mã khuyến mãi:} Nhập mã mới thay thế mã cũ hoặc xóa mã đã áp dụng (\texttt{UC10}).
        \item \textbf{A3 -- Gửi trùng yêu cầu thanh toán:} Khi nhận cùng \texttt{Idempotency-Key}, trả về kết quả giao dịch đầu tiên, không trừ tiền hay cộng điểm lần hai.
        \item \textbf{A4 -- Thanh toán lại đơn \texttt{PENDING}:} Khách có thể thanh toán lại trên đơn hàng chưa hoàn tất (\texttt{UC11a}).
    \end{itemize}
    \item \textbf{Luồng ngoại lệ (Exception Flows):}
    \begin{itemize}
        \item \textbf{E1 -- Chưa đăng nhập / hết phiên:} Chuyển hướng sang trang đăng nhập (\texttt{UC08}), sau khi đăng nhập quay lại trang thanh toán và giữ nguyên giỏ hàng.
        \item \textbf{E2 -- Giỏ hàng rỗng:} Chặn thao tác đặt đơn và thông báo giỏ hàng trống.
        \item \textbf{E3 -- Mã khuyến mãi không hợp lệ:} Thông báo lỗi, giữ nguyên tổng tiền và cho phép tiếp tục thanh toán.
        \item \textbf{E4 -- Không đủ tồn kho:} Rollback Transaction, thông báo món không đủ hàng, không tạo đơn và không thay đổi tồn kho.
        \item \textbf{E5 -- Thiếu \texttt{Idempotency-Key}:} Trả về mã lỗi \texttt{400 Bad Request}, từ chối tạo giao dịch.
        \item \textbf{E6 -- Đơn đã thanh toán:} Chặn thanh toán lại nếu đơn đã ở trạng thái \texttt{PAID}.
        \item \textbf{E7 -- Thanh toán thất bại (\texttt{PAYMENT\_FAILED}):} Giữ đơn ở trạng thái \texttt{PENDING}, giữ nguyên giỏ hàng, không cộng điểm và cấp khóa mới để thử lại.
        \item \textbf{E8 -- Cổng thanh toán không phản hồi:} Giữ đơn ở \texttt{PENDING}, không tự động tạo giao dịch trùng lặp.
    \end{itemize}
\end{itemize}